% Pursuing Stacks (cotes 134-2 to 134-8): archive cote 134-6, site page 112 (PDF page 113), typescript page 430.
% typescript. Draft compared: ps5.tex:7342-7405.
% AI-transcribed by Claude Opus 5.5 (claude-opus-5-5, adaptive thinking, effort high), Message Batches API, batch ps3 (msgbatch_01XdFmuVVBAcgc7EFxdtCqsj), 2026-11-12; not yet checked by eye.
\scanpage{134-6/112}
\mspage{430}

pairing (111) $R\underline{Hom}_k$ . T\hand{h}e point \xout{\unsure{is that}} here is that, whereas\inserted{a}reaso-%
nable $R\underline{Hom}$ can be defined indeed, it \uline{cannot} be computed in terms of
a projective resolution of the first argument as in (113); or equivalent% NOTE: right edge of the scan cuts the line end; no hyphen visible
ly, that for projective $L$ it is not necessarily true that
\[ \underline{Ext}^i_k(L,F) = 0 \quad \text{for } i \hand{>} 0 \ ; \]
this in turn relates to the observation that, contrarily to what
happens for the notion of injective sheaves of modules (on an arbitrary
topos), \uline{it is not true that the property for a sheaf of modules to be
projective is stable under localization}\xout{\unsure{,}} (even for constant sheaf of
rings $k$ on a topos $A^{\wedge}$). Indeed, the localization of $k^{(a)}$ with respect
to $A_{/b}$ (with $a$ and $b$ in $A$) is \xout{$k^{(axb)}$} $k^{(a')}$ where $a'$ is $axb$ viewed
as an object in $A_{/b}$ , and \xout{the} for any sheaf of $k$-modules $F$ on $A_{/b}$
we have
\[ \mathrm{Ext}^i_{\hand{A_{/a}}}(k^{(a')},F) = H^i(A_{/a'=axb} \ , F) \quad , \]
which need not be zero \xout{if} for $i \hand{>} 0$. If it was, this would \xout{mean} imply
that $axb$ is $k$-acyclic (rather, that it's connected components are),
a rather strong property indeed when $axb$ is not in $A$ ...

Thus, when $A$ is not stable under binary products, it does'nt
seem that there exists a pairing (111) as I expected, \xout{unless}\inserted[typed]{except} (possibly)
\inserted{when} there are enough injectives in $A^{\wedge}_{\hand{\mathcal{M}}}$ - a case I do not wish to examine
for the time being, as I am mainly interested\inserted{now} in a formalism using pro-%
jective resolutions instead of injective ones. Anyhow for the purpose
of subsuming the cohomology functor $R\Gamma_A$ under the $R\mathrm{Hom}_k$ formalism,
by formula
% NOTE: in (114) the typed underline under "Hom" of RHom_k is struck through by hand
\[ (114) \qquad R\Gamma_A(F^{\bullet}) \hand{=} R\mathrm{Hom}_k(k_A \ , F^{\bullet}) \hand{\simeq} \mathrm{Hom}_k^{\bullet\bullet}(L^A_{\bullet} \ , F^{\bullet}) \quad , \]
where
\[ (115) \qquad L^A_{\bullet} \hand{\to} k_A \]
is a projective resolution of $k_A$ , it is the pairing $R\mathrm{Hom}_k$ and not
$R\underline{Hom}_k$ which is the \unsure{melevant|relevant} one. Let's recall that a projective reso-%
lution (115) is called a \uline{cointegrator} (for the category $A$, with coeffi% NOTE: right edge of the scan cuts the line end; no hyphen visible
cients in $k$), as by formula (114) it allows indeed\inserted{to express} ``cointegration''
of any \hand{$\mathcal{M}$}-valued presheaf o\over{f}{r} complex of such presheaves (with degrees
bounded from below).

Thus for the time being we just got the pairing $R\mathrm{Hom}_k$ in (110),
and the corresponding functor
\[ \hand{(116)} \qquad \hand{F^{\cdot} \mapsto} R\Gamma_A\hand{(F^{\cdot}) : D^+(A^{\wedge}_{\mathcal{M}}) \longrightarrow D^+(\mathcal{M}) ,} \]

\margin{112}
